\documentclass[aps,onecolumn,nofootinbib,pra]{article}

\usepackage{../../spec_files/arxiv}
\input{../../spec_files/standard_preamble.tex}

\usepackage{import}
\subimport{../../macros/}{all-macros}

\usepackage{minted}

\newcommand{\gausscore}{g}
\newcommand{\gausscoreof}[1]{\gausscore_{#1}}
\newcommand{\gausscoreat}[1]{\gausscore\left[#1\right]}
\newcommand{\gausscoreofat}[2]{\gausscoreof{#1}\left[#2\right]}

\newcommand{\slopevec}{\tilde{\theta}}

\begin{document}
    \title{Squares risk}

    \maketitle
    \date{\today}

    We reformulate results in \cite{goessmann_uniform_2021} based on the formalism in \cite{goessmann_tensor_2026}.


    \section{Quadratic exponential models}

    We consider distributions of variables $\shortcatvariables$ and $Y$ indexed by vector parameters $\slopevec[L]\in\rr^\catorder$ with $L$
    \begin{align*}
        \probofat{\slopevec}{\indexedshortcatvariables,Y=y}
        =
        \frac{1}{\partitionfunction(\slopevec)} \cdot
        \expof{-(\contraction{\slopevec[L],x[L]}-y)^2} \, .
    \end{align*}
    Here we treat the index tuple $\shortcatindices$ as a vector $x[L]$.

    \begin{example}[Gaussian processes]
        When $\shortcatvariables$ and $Y$ are continuous real variables, and the contraction is defined as Lebesque integral, then $\probof{\slopevec}$ is a Gaussian process.


        Here we allow for arbitrary state sets to $\shortcatvariables$ and $Y$, but are specially interested in discrete cases (where the contraction is the sum).
    \end{example}

    \subsection{Cross entropy}

    The cross entropy between an empirical distribution $\empdistribution$ to the data $\datamap=\dataset$ and the model $\probof{\slopevec}$ is
    \begin{align*}
        \centropyof{\empdistribution}{\probof{\slopevec}}
        &= \frac{1}{\datdim} \sum_{\datindexin}
        -\lnof{\probofat{\slopevec}{\indexedshortcatvariables^\datindex,Y=y^\datindex}} \\
        &= \lnof{\partitionfunction(\slopevec)}
        - \frac{1}{\datdim} \sum_{\datindexin}
        (\contraction{\slopevec[L],\catindex^\datindex} - y^\datindex)^2 \, .
    \end{align*}

    The cross entropy differs by the term $\lnof{\partitionfunction(\slopevec)}$ from the squares risk
    \begin{align*}
        \frac{1}{\datdim} \sum_{\datindexin}
        \left(\contraction{\slopevec[L],\catindex^\datindex} - y^\datindex\right)^2 \, .
    \end{align*}

    \subsection{Conditioned likelihood}

    When $Y$ is real valued, then the distribution of $Y$ conditioned on $\indexedshortcatvariables$ is
    \begin{align*}
        \probofat{\slopevec}{Y=y|\indexedshortcatvariables} \, ,
    \end{align*}
    that is a normal distribution with $\sigma=\sqrt{\frac{1}{2}}$.

    The log likelihood of the data $\datamap$ conditioned on the variables $\shortcatvariables$ is
    \begin{align*}
        \lnof{\sqrt{\pi}} +
        \frac{1}{\datdim} \sum_{\datindexin}
        \left(\contraction{\slopevec[L],\catindex^\datindex} - y^\datindex\right)^2 \, .
    \end{align*}
    Thus, the difference to the squares risk is a constant term $\lnof{\sqrt{\pi}}$, which can be ignored when optimizing over $\slopevec$.

    \subsection{Exponential family}

    By factorization of the squares term we get
    \begin{align*}
        \probofat{\slopevec}{\indexedshortcatvariables,Y=y}
        =
        \frac{1}{\partitionfunction(\slopevec)}
        \expof{
            -\contraction{\slopevec[L],x[L]}^2
            + 2 y \contraction{\slopevec[L],x[L]}
            - y^2
        } \, .
    \end{align*}
    This is a member of the exponential family with the statistic
    \begin{align*}
        \sstat(x,y)
        = (-x[L]\otimes x[\tilde[L]], 2y \cdot x[L], -y^2)
    \end{align*}
    and the canonical parameter
    \begin{align*}
        \canparam(\slopevec)
        = (\slopevec[L]\otimes\slopevec[\tilde{L}], \slopevec[L], 1) \, .
    \end{align*}


    \section{Inference}

    The special case of Gaussian processes enable training methods like alternating least squares, see \cite{goessmann_uniform_2021}.
    They require sparsity assumptions by restricting $\slopevec$ to admit tensor network decompositions.

    \section{Probabilistic guarantees}

    From the basic inequality
    \begin{align*}
        \centropyof{\empdistribution}{\probof{\hat{\slopevec}}}
        \leq
        \centropyof{\empdistribution}{\probof{\slopevec^*}}
    \end{align*}
    we derive an inequality consisting product and multiplier processes, scaling with $\|\hat{\slopevec}-\slopevec^*\|^2$ and $\|\hat{\slopevec}-\slopevec^*\|$.
    In combination with uniform concentration inequalities (involving e.g. Mendelsons Lambda functional), we then get a bound on $\|\hat{\slopevec}-\slopevec^*\|$.



    \bibliographystyle{plainnat}
    \bibliography{../../references.bib}


\end{document}